% GENERATED FILE. Edit canonical lessons, then run npm run generate:books.
% canonical-source-hash: fnv1a64:2fc424c6b76501f9
% canonical-lessons: IT-C521-pittura, IT-C521-pennello, IT-C521-tela, IT-C521-cavalletto, IT-C521-acquerello

\chapter{Pennelli e colori}
\label{ch:it-pennelli-e-colori}
\hlchaptermodality{521}

\begin{chapteropening}
\textbf{By the end of this chapter:} \emph{I can talk about painting as a hobby.}

Complete the last lesson of chapter 521: I can talk about painting as a hobby.
\end{chapteropening}

\section[la pittura]{la pittura — painting; paint}

\label{lesson:IT-C521-pittura}



\begin{quote}
Before the new one: say the Italian for to make one's debut, then the Italian for director.
\end{quote}

\subsection*{You'll want to know: la pittura}
\textbf{la pittura} — \textquotedblleft{}painting; paint\textquotedblright{}.

\begin{quote}
La pittura mi rilassa molto.
\end{quote}

\begin{quote}
\emph{Painting relaxes me a lot.}
\end{quote}

Also wet paint: \emph{pittura fresca}.

\begin{cousinweb}
\emph{pittura} is inherited from Latin \emph{pictūra} \textquotedblleft{}the art of painting\textquotedblright{} (EtymDB), from \emph{pingere} \textquotedblleft{}to paint\textquotedblright{}.
\end{cousinweb}

\begin{grammarlens}[title={What you've built}]
One more word for telling, explaining and describing at B1.
\end{grammarlens}

\subsection*{Your turn}
\begin{itemize}
\raggedright
  \item \emph{Say it:} \emph{la pittura}
  \item \emph{Say it:} \emph{la pittura}, once more
  \item \emph{Say it:} a sentence of your own with \emph{la pittura}
  \item \emph{Recall:} say the Italian for to make one's debut, then the Italian for director, then say \emph{la pittura} again
\end{itemize}

\begin{tcolorbox}[breakable,colback=teal!4,colframe=teal!35!black,title={Before you move on}]
What does la pittura mean? (\textquotedblleft{}Painting; paint\textquotedblright{}.) Say it once more.
\end{tcolorbox}
\section[il pennello]{il pennello — paintbrush}

\label{lesson:IT-C521-pennello}



\begin{quote}
Before the new one: say the Italian for director, then the Italian for painting.
\end{quote}

\subsection*{You'll want to know: il pennello}
\textbf{il pennello} — \textquotedblleft{}paintbrush\textquotedblright{}.

\begin{quote}
Lavo i pennelli con l'acqua calda.
\end{quote}

\begin{quote}
\emph{I wash the brushes in hot water.}
\end{quote}

\emph{A pennello} — perfectly: \emph{ti sta a pennello}.

\begin{cousinweb}
\emph{pennello} is inherited from Vulgar Latin \emph{pēnicellus} \textquotedblleft{}a little brush\textquotedblright{} (Lewis \& Short has \emph{pēnicillus}), a little form of a word for a tail.
\end{cousinweb}

\begin{grammarlens}[title={What you've built}]
One more word for telling, explaining and describing at B1.
\end{grammarlens}

\subsection*{Your turn}
\begin{itemize}
\raggedright
  \item \emph{Say it:} \emph{il pennello}
  \item \emph{Say it:} \emph{il pennello}, once more
  \item \emph{Say it:} a sentence of your own with \emph{il pennello}
  \item \emph{Recall:} say the Italian for director, then the Italian for painting, then say \emph{il pennello} again
\end{itemize}

\begin{tcolorbox}[breakable,colback=teal!4,colframe=teal!35!black,title={Before you move on}]
What does il pennello mean? (\textquotedblleft{}Paintbrush\textquotedblright{}.) Say it once more.
\end{tcolorbox}
\section[la tela]{la tela — canvas; cloth}

\label{lesson:IT-C521-tela}



\begin{quote}
Before the new one: say the Italian for painting, then the Italian for paintbrush.
\end{quote}

\subsection*{You'll want to know: la tela}
\textbf{la tela} — \textquotedblleft{}canvas; cloth\textquotedblright{}.

\begin{quote}
Ho dipinto un paesaggio su una tela grande.
\end{quote}

\begin{quote}
\emph{I painted a landscape on a large canvas.}
\end{quote}

Also cloth: \emph{una borsa di tela}.

\begin{cousinweb}
\emph{tela} is inherited from Latin \emph{tēla} \textquotedblleft{}a web, a loom\textquotedblright{} (EtymDB), from \emph{texere} \textquotedblleft{}to weave\textquotedblright{}.
\end{cousinweb}

\begin{grammarlens}[title={What you've built}]
One more word for telling, explaining and describing at B1.
\end{grammarlens}

\subsection*{Your turn}
\begin{itemize}
\raggedright
  \item \emph{Say it:} \emph{la tela}
  \item \emph{Say it:} \emph{la tela}, once more
  \item \emph{Say it:} a sentence of your own with \emph{la tela}
  \item \emph{Recall:} say the Italian for painting, then the Italian for paintbrush, then say \emph{la tela} again
\end{itemize}

\begin{tcolorbox}[breakable,colback=teal!4,colframe=teal!35!black,title={Before you move on}]
What does la tela mean? (\textquotedblleft{}Canvas; cloth\textquotedblright{}.) Say it once more.
\end{tcolorbox}
\section[il cavalletto]{il cavalletto — easel; stand}

\label{lesson:IT-C521-cavalletto}



\begin{quote}
Before the new one: say the Italian for paintbrush, then the Italian for canvas.
\end{quote}

\subsection*{You'll want to know: il cavalletto}
\textbf{il cavalletto} — \textquotedblleft{}easel; stand\textquotedblright{}.

\begin{quote}
Metto la tela sul cavalletto.
\end{quote}

\begin{quote}
\emph{I put the canvas on the easel.}
\end{quote}

A little \emph{cavallo} \textquotedblleft{}horse\textquotedblright{}; any frame on legs.

\begin{cousinweb}
\emph{cavalletto} is a little \emph{cavallo} \textquotedblleft{}horse\textquotedblright{}, inherited from Latin \emph{caballus} \textquotedblleft{}a nag\textquotedblright{} (Lewis \& Short): a frame on legs.
\end{cousinweb}

\begin{grammarlens}[title={What you've built}]
One more word for telling, explaining and describing at B1.
\end{grammarlens}

\subsection*{Your turn}
\begin{itemize}
\raggedright
  \item \emph{Say it:} \emph{il cavalletto}
  \item \emph{Say it:} \emph{il cavalletto}, once more
  \item \emph{Say it:} a sentence of your own with \emph{il cavalletto}
  \item \emph{Recall:} say the Italian for paintbrush, then the Italian for canvas, then say \emph{il cavalletto} again
\end{itemize}

\begin{tcolorbox}[breakable,colback=teal!4,colframe=teal!35!black,title={Before you move on}]
What does il cavalletto mean? (\textquotedblleft{}Easel; stand\textquotedblright{}.) Say it once more.
\end{tcolorbox}
\section[l'acquerello]{l'acquerello — watercolour}

\label{lesson:IT-C521-acquerello}



\begin{quote}
Before the new one: say the Italian for canvas, then the Italian for easel.
\end{quote}

\subsection*{You'll want to know: l'acquerello}
\textbf{l'acquerello} — \textquotedblleft{}watercolour\textquotedblright{}.

\begin{quote}
Dipingo all'acquerello paesaggi di mare.
\end{quote}

\begin{quote}
\emph{I paint seascapes in watercolour.}
\end{quote}

\emph{All'acquerello} — in watercolour.

\begin{cousinweb}
\emph{acquerello} is built on \emph{acqua} \textquotedblleft{}water\textquotedblright{} (EtymDB), inherited from Latin \emph{aqua}.
\end{cousinweb}

\begin{grammarlens}[title={What you've built}]
One more word for telling, explaining and describing at B1.
\end{grammarlens}

\subsection*{Your turn}
\begin{itemize}
\raggedright
  \item \emph{Say it:} \emph{l'acquerello}
  \item \emph{Say it:} \emph{l'acquerello}, once more
  \item \emph{Say it:} a sentence of your own with \emph{l'acquerello}
  \item \emph{Recall:} say the Italian for canvas, then the Italian for easel, then say \emph{l'acquerello} again
\end{itemize}

\begin{tcolorbox}[breakable,colback=teal!4,colframe=teal!35!black,title={Before you move on}]
What does l'acquerello mean? (\textquotedblleft{}Watercolour\textquotedblright{}.) Say it once more.
\end{tcolorbox}
